\prefacesection{Abstract}

Accurate electricity demand forecasting is critical for the reliable and efficient operation of power grids, particularly in compact, highly urbanised nations such as Singapore where climate conditions remain consistently tropical and demand fluctuations are strongly coupled to weather patterns. Despite the growing availability of high-resolution meteorological data and advances in machine learning, systematic comparative studies that rigorously evaluate the predictive contribution of climatic variables to short-term electricity demand remain limited for this geographical context.

This study presents a comprehensive machine learning pipeline for daily electricity demand forecasting in Singapore, leveraging weather-derived and temporal features extracted from data spanning February 2023 to December 2025. A feature set of 16 variables was constructed, encompassing temperature, humidity, a Comfort Climate Index (CCI), calendar-based indicators (day of week, month, season, public holidays, and weekends), and historical demand lag and rolling statistics. Eight model configurations were trained and evaluated under a chronological 80/20 train-test split: a Lag-1 baseline, Linear Regression, Random Forest, XGBoost, and CatBoost with default settings, plus PPSO-tuned variants of each tree-based model, where every tunable model received an identical Phasor Particle Swarm Optimisation (PPSO) budget to ensure a fair comparison. Models were assessed using Mean Absolute Error (MAE), Root Mean Squared Error (RMSE), Coefficient of Determination ($R^2$), Mean Absolute Percentage Error (MAPE), Relative Absolute Error (RAE), and Willmott Index (WI). The best model is further deployed through an interactive web portal (FastAPI backend and React frontend) that provides historical analysis and forward-looking demand outlooks driven by live weather forecasts.

The PPSO-optimised CatBoost model achieved the strongest predictive performance, attaining an $R^2$ of 0.914, an MAE of 1{,}496.2 MWh, and an RMSE of 1{,}895.5 MWh on the held-out test set, outperforming all other configurations. Under the identical tuning budget, PPSO also improved XGBoost but not Random Forest, showing that the benefit of meta-heuristic tuning is model-dependent. These results demonstrate that automated hyperparameter tuning via PPSO yields meaningful gains in gradient-boosted tree models and that weather-derived features, particularly temperature and its interaction with humidity, are significant predictors of electricity demand in Singapore. The findings have practical implications for grid operators and energy planners seeking data-driven approaches to short-term load forecasting.

\prefacesection{Acknowledgements}

The authors would like to express sincere gratitude to Dr.\ Young Lee for his guidance, constructive feedback, and continued support throughout this project. We also thank the Singapore Energy Market Authority (EMA) for publicly providing the electricity demand data that formed the basis of this study, and the open-source communities behind scikit-learn, XGBoost, CatBoost, and the Open-Meteo weather API, whose tools made this research possible.

\afterpreface